\section{Implementation}
\label{sec:implementation}

\subsection{Stack and layout}

\begin{center}
\small
\begin{tabular}{@{}ll@{}}
\toprule
pygeoapi 0.24.0 & OGC API framework: routing, job management, OpenAPI \\
shapely 2.1     & geometry: union, difference, intersection, area \\
pyproj 3.8      & per-request equal-area projection \\
rasterio        & raster windowed reads and masking \\
TinyDB          & record store and persisted job store \\
gunicorn        & WSGI server, 3 workers, 300\,s timeout \\
\bottomrule
\end{tabular}
\end{center}

\begin{lstlisting}
bhoonidhi-records/
|- pygeoapi/config.yml                 server, Records resource, 6 processes
|- pygeoapi/processes/
|  |- catalogue_spatial.py             4 footprint processes
|  `- raster_zonal.py                  2 raster processes
|- scripts/
|  |- build_catalogue.py               59 collection records, provenance-tagged
|  |- fetch_footprints.py              STAC client (bhoonidhi | open)
|  |- probe_gap.sh                     Method 1: specification audit
|  |- probe_baseline.sh                Method 2: ISRO + peer baseline
|  |- measure_redundancy.py            Method 3: the 92% figure
|  |- serve_local.sh                   start | stop | status  (tested path)
|  |- run_tests.sh                     the 100-test compliance suite
|  `- demo.sh                          the live walkthrough
|- tests/                              records | processes | raster
|- data/                               records.tinydb, footprints.geojson, dem.tif
`- docs/
   |- GAP_ANALYSIS.md                  the long-form gap argument
   |- evidence/                        saved probe responses + measurements
   `- report/                          this document
\end{lstlisting}

\subsection{Process structure}

Each process is a pygeoapi \code{BaseProcessor} subclass declaring its
metadata --- identifier, title, description, input and output schemas,
keywords and a worked example --- and implementing \code{execute()}. pygeoapi
derives the \ep{/processes} listing, the per-process description, the JSON
Schema validation of inputs and the OpenAPI document from that declaration,
which is why conformance comes from the framework rather than from
hand-written routing.

A representative shape, from \code{minimal-scene-set}:

\begin{lstlisting}
def execute(self, data, outputs=None):
    aoi      = _aoi(data.get("aoi"))            # bbox | Polygon | Feature
    target   = float(data.get("target_fraction") or 0.99)
    fps      = _load(coll, start, end, aoi)     # footprints intersecting AOI
    to_m, from_m = _equal_area(aoi)             # per-request LAEA
    aoi_m    = shp_transform(to_m, aoi)

    # greedy set cover, in metres
    while remaining and len(chosen) < cap:
        if covered is not None and covered.area / aoi_m.area >= target:
            break
        best, best_gain = None, 0.0
        for c in remaining:                     # largest new area wins
            gain = c["g"].area if covered is None \
                   else c["g"].difference(covered).area
            if gain > best_gain:
                best, best_gain = c, gain
        if best is None or best_gain < 1_000:   # under 1000 m2 is noise
            break
        covered = unary_union([covered, best["g"]]) if covered else best["g"]
        chosen.append({... "new_area_km2": ..., "cumulative_coverage_fraction": ...})

    return "application/json", { ..., "algorithm": "greedy set cover; ...",
                                 "provenance": _provenance(...) }
\end{lstlisting}

Three details worth drawing out. Input coercion (\code{\_aoi}) accepts a bare
bbox array, a GeoJSON \code{Polygon} or a GeoJSON \code{Feature}, because
clients reasonably send all three, and tests cover each. Invalid geometry from
upstream is repaired with \code{buffer(0)} rather than crashing the request.
And the whole optimisation runs in projected metres, never in degrees.

\subsection{Job management and asynchronous execution}

OGC API --- Processes requires more than a synchronous endpoint. The service
declares \code{jobControlOptions: [sync-execute, async-execute]} on every
process and implements the full lifecycle:

\begin{center}
\small
\begin{tabularx}{\textwidth}{@{}l>{\raggedright\arraybackslash}X@{}}
\toprule
\textbf{Request} & \textbf{Behaviour} \\
\midrule
\code{POST .../execution}
  & synchronous: HTTP 200 with the result document. \\
\code{POST .../execution}\newline with \code{Prefer: respond-async}
  & HTTP \textbf{201 Created}, a \code{Location} header pointing at the job,
    \code{Preference-Applied: respond-async}, and a body carrying
    \code{jobID} and \code{status: accepted}. \\
\code{GET /jobs/\{id\}}
  & job status document: \code{status}, \code{progress}, \code{created},
    \code{started}, \code{finished}, and links. \\
\code{GET /jobs/\{id\}/results}
  & the result document once \code{status: successful}. \\
\bottomrule
\end{tabularx}
\end{center}

Jobs are persisted through pygeoapi's TinyDB manager, so asynchronous
execution survives a restart and \ep{/jobs/\{id\}} keeps answering. A
compliance test walks the entire lifecycle --- submit async, poll to
\code{successful}, fetch \ep{/results} --- rather than merely checking that a
201 came back.

\subsection{Defects found during development, and fixed}
\label{sec:defects}

Three real defects were found and each is now guarded by a test. They are
recorded because they are the part of the work that was actually difficult, and
because a compliance suite that never caught anything is not evidence of much.

\paragraph{1. Every process execution returned HTTP 400 after a reboot.}
The most serious of the three, and invisible in development. pygeoapi's TinyDB
job manager writes each job's result into \code{output\_dir} and \textbf{does
not create that directory}. The configured path was under \code{/tmp}, which
the operating system clears on reboot. The directory existed during the
session it was first created in, so everything passed; after a reboot every
single \code{POST .../execution} --- synchronous and asynchronous alike ---
failed with an opaque message:

\begin{lstlisting}
HTTP 400 {"code":"InvalidParameterValue","description":"Error executing
process: [Errno 2] No such file or directory:
'/tmp/bhoonidhi-records-output/zonal-statistics-0ee588ce-...'"}
\end{lstlisting}

This took 39 of the 100 tests down at once. The fix moves the job store under
\code{runtime/jobs} via a \code{\$\{JOBS\_DIR\}} substitution in
\code{config.yml}, and \code{serve\_local.sh} now creates it before starting
gunicorn --- so the directory is a startup precondition rather than an
accident of history, and job records survive a reboot as a side benefit.

\paragraph{2. A hairline sliver was reported as a gap while also reporting
\code{fully\_covered: true}.} Differencing two near-identical polygons leaves
slivers with effectively zero area. \code{aoi-coverage} would return
\code{fully\_covered: true} and simultaneously hand back a gap geometry ---
self-contradictory output that any client would mishandle. Anything under
0.01\,km\textsuperscript{2} (one hectare) is now collapsed to no gap at all,
the threshold is published in the response as \code{gap\_threshold\_km2} so
the behaviour is inspectable rather than magic, and a test asserts that
\code{fully\_covered} never coexists with a gap geometry.

\paragraph{3. The zero-footprint branch returned a different JSON shape.}
When no footprints matched the filters, \code{aoi-coverage} returned an object
with different keys from the populated case --- which would break any client
the first time a search came back empty. Both branches now return the same
schema, with the empty case carrying \code{scene\_count: 0},
\code{gap\_fraction: 1.0}, the AOI itself as the gap, and an explanatory
\code{message}. A test asserts shape equality across both branches.

\subsection{Deployment}

\begin{lstlisting}
make install      # Python venv
make data         # build the 59 catalogue records
make footprints   # harvest real acquisition footprints
make up           # start the service on :5003
make test         # 100 compliance tests
make demo         # the whole walkthrough
\end{lstlisting}

\code{scripts/serve\_local.sh} is the tested path. It regenerates the OpenAPI
document from \code{config.yml} on every start, so the published API
description cannot drift from the configuration, and it detaches gunicorn with
\code{setsid} and a closed stdin so that \code{serve\_local.sh start | tail}
does not hang on a pipe the server keeps alive.

For a publicly reachable URL, \code{scripts/expose\_public.sh} opens a
Cloudflare quick tunnel. The resulting hostname is ephemeral --- a fresh one
per run --- which is noted in Section~\ref{sec:limitations}.
